\scanpage{130}
\uline{Corollaire} Soit $L$ un ss-corps de $\overline Q$. Pour qu'on ait $L\subset\mathbb L$, il f. et s. que $L$ soit tot réel, ou extension quadratique purement imaginaire d'un corps tot. réel.

Suffisance triviale par (iii). Nécessité: Par théorème de Galois, $L$ est une extension de degré $1$ ou $2$ du corps tot. réel $K=L\cap\mathbb K$, et il reste à prouver que lorsque $L$ en est une extension quadratique, cette dernière est \unsure{tot\textsuperscript{t}|tot.} \uline{imaginaire}.
\margin{\diagram{\begin{tikzcd}[row sep=small,column sep=small]
\mathbb K \arrow[r,dash] \arrow[d,dash] & \mathbb K(L) \arrow[r,dash] \arrow[d,dash] & \mathbb L \\
K=\mathbb K\cap\unsure{L|\mathbb L} \arrow[r,dash] \arrow[d,dash] & L \\
Q
\end{tikzcd}}}
\deleted{\unsure{On a}} Or \deleted{\unsure{on voit aisément que}} on est ramené à prouver
\deleted{\uline{Lemme} Soit $L$ une extension quadratique d'un corps tot. réel $K$ de $\mathbb L$}
(quitte à remplacer $L$ par un conjugué, \inserted{ce qui ne dérange pas $K$,}) que $L$ est imaginaire. Or s'il était réel, comme $\mathbb K$ l'est, $\mathbb K(L)$ le serait, or $\mathbb K(L)=\mathbb L$, absurde.

\uline{NB} $\mathbb L$ est la plus grande extension galoisienne de $Q$, telle que dans son groupe de Galois $G'$, l'image de $\tau$ (conjugaison complexe) soit \uline{centrale}. Donc $\mathbb L$ contient l'extension abélienne maximale de $Q$, i.e. le corps des racines de l'unité.
\mspage{129}
